% SPDX-FileCopyrightText: 2026 Bart van de Lint
% SPDX-License-Identifier: AGPL-3.0-or-later
\section{Introduction}
\label{sec:introduction}

Shipping burns more than 300 million tonnes of fuel a year and emitted \num{1056} million tonnes of CO$_2$ in 2018 \citep{faber2020}. The International Maritime Organization aims for net-zero greenhouse gas emissions from international shipping by or around 2050 \citep{imo2023}. Fuels without fossil carbon, such as e-ammonia, made from renewable hydrogen and nitrogen from the air, can do this, but they cost three to four times as much per unit of energy as fuel oil (\cref{sec:costs}). Wind can pay part of that bill. Flettner rotors, tall spinning cylinders that turn the wind into a forward force, and rigid or semi-rigid wing sails are installed on cargo ships, and towing kites have flown on several \citep{emsa2023}. Published reductions in CO$_2$ emissions range from 2 to 48\,\% \citep{huang2025}, because the studies differ in ship, route and method and report their results inconsistently.

The rigs differ in what they give and what they cost. A wing sail and a rotor stand on deck. They take deck space and cargo-handling access, and they pull at a height of tens of metres in the apparent wind, the wind felt on board, which the ship's own speed turns forward. A rotor makes about five times the lift, the force across the wind, of a wing of the same area, but needs power to spin. A kite flies hundreds of metres up on a tether and sweeps across the wind in figures of eight, which multiply the wind it feels and so its pull. It takes almost no deck space but has to be launched and recovered, and its textile canopy wears out in flight. Each rig also pushes the ship sideways and heels it, leaning it over. To balance the side force the hull sails at a small angle to its course, the leeway angle, and the rudder holds the course. Both add drag, which costs a share of the rig's gain \citep{kramer2022}.

We compare the three rigs on cost and fuel energy across ships, routes, fuels and assumptions, with one model that takes a ship design to its cost per tonne-kilometre. It covers the hull and engine, the forces of each rig and their balance by hull and rudder, weather routing (the choice of path and speed through the weather) and the full voyage cost of capital, crew, ports, fuel, kite wear and cargo interest. The ships sail three years of the ERA5 reconstruction of past weather on six Atlantic routes. There are thirteen standard ship classes, from an \SI{85}{m} coaster to a Megamax container ship of \num{24000} twenty-foot containers (TEU). Each has a mechanical drive on very low sulphur fuel oil (VLSFO), today's standard fuel, and on e-ammonia, and every rig is sized for every class, route and fuel.

Three choices shape the answer. First, every input carries its source or says that it has none, and all inputs are collected in one table (\cref{sec:inputs}). Second, the kite is modelled quasi-statically and assumed to hold its design shape. A vortex method gives its aerodynamics, and its forces balance at every point of a steered flight pattern, from figures of eight to slow sweeps along the edge of the wind window. A search optimises the pattern for every wind and course, and the model is calibrated to measured flight. Third, cost is counted per tonne-kilometre, with time costed as capital, crew and cargo interest. The router therefore trades fuel against speed, and a rig that saves fuel also buys time. The terms of both fields are explained where they first appear and collected in the \hyperref[sec:nomenclature]{Nomenclature}.

The result is stated as a condition: what each rig saves if it reaches goals for life, canopy price, reliability and launch that development can still change (\cref{sec:res_goals}). The paper contributes four things. It puts all three rigs on one model and cost basis, with each rig sized close to its own optimum. It maps the cost saving of each rig, the rig that comes out ahead and the rig sizes over thirteen classes, six routes and two fuels (\cref{sec:res_fuel,sec:res_size}), and shows in one section that a VTOL kite, a rigid wing that takes off and lands vertically on its own propellers, is too expensive at ship scale (\cref{sec:res_rigid}). It sweeps the assumptions on which the comparison turns and gives the point at which each rig stops paying if it misses a goal (\cref{sec:res_situations,sec:res_missed}). And a companion website at \tbd{URL} lets readers change the assumptions themselves.

\Cref{sec:related} reviews earlier comparisons. \Cref{sec:model} describes the model, its costs and the goals. \Cref{sec:results} says which rig wins where and why, \cref{sec:certainty} how sure that is, and \cref{sec:discussion} gives the limits of the model. \Cref{sec:industry} draws out what this means for shipowners and makers, and \cref{sec:conclusions} concludes. The appendices give the formulas and calibrations behind the method (\cref{sec:details}), the detail behind the limits (\cref{sec:disc_details}), further results (\cref{sec:further}), the model checks (\cref{sec:checks}) and the inputs (\cref{sec:inputs}).
